\begin{table}[htbp]
\centering
\begin{tabular}{rrrrrrrr}
\toprule
& \multicolumn{3}{c}{reference path (s)} & \multicolumn{3}{c}{incremental path (s)} & \\
\cmidrule(lr){2-4}\cmidrule(lr){5-7}
$n$ & median & max & per event & median & max & per event & ratio \\
\midrule
3 & 0.010 & 0.010 & 0.0097 & 0.012 & 0.013 & 0.0124 & 1.28 \\
4 & 0.105 & 0.165 & 0.0154 & 0.148 & 0.176 & 0.0216 & 1.40 \\
5 & 0.529 & 0.775 & 0.0279 & 0.798 & 1.092 & 0.0394 & 1.51 \\
6 & 2.286 & 3.335 & 0.0548 & 3.622 & 4.776 & 0.0785 & 1.58 \\
7 & 7.967 & 12.791 & 0.1077 & 14.619 & 19.779 & 0.1783 & 1.83 \\
8 & 28.239 & 32.618 & 0.2097 & 41.563 & 52.905 & 0.3012 & 1.47 \\
\midrule
\multicolumn{8}{l}{Power-law fit of median time: reference $t \propto n^{8.0}$, incremental $t \propto n^{8.3}$.} \\
\bottomrule
\end{tabular}
\caption{Wall-clock time, single-threaded. ``per event'' is median total time divided by the number of events. ``ratio'' is incremental median over reference median: values above $1$ mean the incremental path is slower, because it also recomputes the slice to assert agreement.}
\label{tab:timing}
\end{table}
